\documentclass[11pt,twocolumn]{article}

%% ─── Packages ────────────────────────────────────────────────────────────────
\usepackage[margin=1in,columnsep=0.25in]{geometry}
\usepackage{amsmath,amssymb,amsthm}
\usepackage{algorithm}
\usepackage{algpseudocode}
\usepackage{booktabs}
\usepackage{tabularx}
\usepackage{hyperref}
\usepackage{xcolor}
\usepackage{microtype}
\usepackage{enumitem}
\usepackage{cite}
\usepackage{graphicx}
\usepackage{subcaption}
\usepackage{setspace}
\usepackage{abstract}

%% ─── Theorem environments ────────────────────────────────────────────────────
\newtheorem{definition}{Definition}
\newtheorem{proposition}{Proposition}

%% ─── Metadata ────────────────────────────────────────────────────────────────
\title{\textbf{Chainmail: A Unified Open-Source Workbench for Software Supply
Chain Security}\\
\large Novel Multi-Dimensional Risk Scoring, Graph-Theoretic License Conflict
Propagation, and Automated VEX Generation}

\author{
  Sai Sravan Cherukuri\\
  Independent Researcher\\
  \texttt{saisravan@gmail.com}
}

\date{April 2026 \quad arXiv preprint \texttt{cs.CR / cs.SE}}

%% ─────────────────────────────────────────────────────────────────────────────
\begin{document}
\maketitle

%% ─── ABSTRACT ────────────────────────────────────────────────────────────────
\begin{abstract}
Software supply chain attacks have grown explosively, with targeted incidents
increasing by 742\% between 2019 and 2022. Regulatory mandates now require SBOMs,
provenance attestation, and VEX statements from every team shipping software to
federal agencies, yet open-source maintainers still lack unified, cost-free
tooling to meet these demands. I built \textbf{Chainmail} to close that gap: a
single open-source workbench that brings eight supply chain security capabilities
together in one place.

This paper describes three technical contributions at Chainmail's core.
First, I introduce a formal \emph{Multi-Dimensional Supply Chain Risk Score
(SCRS)} that distills five orthogonal risk dimensions into one normalized A-to-F
grade, with proven monotonicity and boundedness properties. Second, I present a
\emph{graph-theoretic License Conflict Propagation (LCP)} framework that reframes
FOSS license compatibility as a reachability problem on directed acyclic dependency
graphs, catching transitive conflicts that flat-list scanners miss entirely. Third,
I describe an \emph{Automated VEX Generation} algorithm that maps CVE records
against versioned SBOM components using PURL-based version-range intersection,
producing CSAF~2.0-compliant statements in under two seconds.

I also present a quantitative 2026 Mandate Readiness Score benchmarked against
24 CISA-defined criteria. Evaluation across the top 500 npm packages shows
Chainmail finds 23\% more license conflicts than existing tools and produces risk
grades correlating at Spearman $\rho = 0.81$ with independently assessed security
scores. A live implementation is available at
\url{https://chainmail.dev}.
To my knowledge, no other open-source tool unifies all eight capabilities---SBOM
health, license conflict detection, risk scoring, VEX generation, provenance setup,
SBOM diffing, license inheritance visualization, and mandate readiness---in a
single API-driven workbench.
\end{abstract}

\noindent\textbf{Keywords:} software supply chain security, SBOM, VEX,
license compliance, risk scoring, dependency graph analysis, open-source
security, SLSA provenance, Executive Order 14028, CISA compliance

\bigskip\hrule\bigskip

%% ─── SECTION 1 ───────────────────────────────────────────────────────────────
\section{Introduction}

The software supply chain has emerged as one of the most consequential attack
surfaces in modern computing. The SolarWinds incident of 2020, the Log4Shell
vulnerability of 2021, and the XZ Utils backdoor of 2024 each demonstrated that
compromising a single upstream dependency can cascade through thousands of
downstream consumers with devastating effect~\cite{peisert2021,chen2022log4shell,binsalleeh2024xz}.
The Open Source Security Foundation (OpenSSF) estimates that 96\% of modern
enterprise applications contain open-source code, and the average project depends
on 203 third-party packages across all transitive levels~\cite{openssf2024}.

Governments have responded with regulatory mandates. U.S. Executive Order
14028~\cite{eo14028} requires Software Bills of Materials (SBOMs) for all
software sold to federal agencies. CISA published minimum SBOM requirements in
2023~\cite{cisa2023sbom}, and the EU Cyber Resilience Act (CRA) mandates similar
attestation requirements for software sold in European markets, taking effect in
2027~\cite{eu2024cra}. Yet despite this regulatory pressure, the ecosystem of
open-source supply chain security tooling remains fragmented, single-purpose, and
frequently behind enterprise paywalls.

Commercial tools---Snyk, FOSSA, Sonatype Nexus IQ, and JFrog Xray---provide
capable coverage but at licensing costs that exclude independent maintainers and
small teams. Open-source alternatives are siloed: Syft generates SBOMs; Grype
scans vulnerabilities; ScanCode detects licenses; SLSA tooling handles
provenance. No single open-source tool integrates these capabilities, computes a
holistic risk score, auto-generates VEX statements, or tracks compliance
readiness against 2026 federal mandates.

I built \textbf{Chainmail} to address this gap directly. It is a single
open-source workbench providing eight security capabilities through one REST API
and web interface. Three of those capabilities rest on novel technical
contributions that this paper documents in detail:

\begin{enumerate}
  \item \textbf{Supply Chain Risk Score (SCRS) Model:} A formal weighted
    composite function over five orthogonal risk dimensions yielding a normalized
    score $\in [0,100]$ mapped to an A--F letter grade via empirically derived
    thresholds.
  \item \textbf{License Conflict Propagation (LCP) Framework:} A
    graph-theoretic formalization of FOSS license compatibility as a reachability
    problem on directed acyclic dependency graphs, with a polynomial-time conflict
    detection algorithm.
  \item \textbf{Automated VEX Generation Algorithm:} An algorithm for
    automatically mapping CVE records against versioned SBOM components using
    PURL-based version-range intersection, producing CISA CSAF~2.0-compliant VEX
    statements without manual analyst intervention.
\end{enumerate}

Section~2 surveys related work. Section~3 describes the system architecture.
Sections~4--6 develop the three core technical contributions. Section~7 covers the
Readiness Score. Section~8 reports experimental evaluation, Section~9 discusses
limitations and future directions, and Section~10 concludes.

%% ─── SECTION 2 ───────────────────────────────────────────────────────────────
\section{Background and Related Work}

\subsection{Software Bill of Materials}

A Software Bill of Materials is a structured, machine-readable inventory of all
software components and their relationships within a given artifact. Two formats
dominate: CycloneDX (OWASP)~\cite{cyclonedx} and SPDX (Linux
Foundation)~\cite{spdx}. Tools such as Syft~\cite{syft}, cdxgen~\cite{cdxgen},
and Trivy~\cite{trivy} generate SBOMs from container images, source trees, and
package lock files. While generation is increasingly mature, the tooling for
\emph{analyzing} SBOMs holistically remains nascent. Chainmail ingests both
CycloneDX and SPDX formats and provides the first open-source multi-dimensional
analysis pipeline.

\subsection{Vulnerability Management and VEX}

CVE records are the primary artifact for tracking known
vulnerabilities~\cite{mitre2024cve}. CVSS provides severity
scores~\cite{cvss}. The Vulnerability Exploitability eXchange (VEX) format,
standardized by CISA via CSAF~2.0~\cite{csaf20}, allows vendors to communicate
whether a known vulnerability is actually exploitable in their specific build
context. VEX generation has historically been a manual, analyst-driven process.
Chainmail's automated VEX algorithm is, to my knowledge, the first open-source
implementation of systematic VEX generation from SBOM data.

\subsection{License Compliance}

FOSS license compliance has been studied extensively in the legal
literature~\cite{rosen2005}, but formal computational models remain limited.
Kapitsaki et al.~\cite{kapitsaki2017} model license compatibility as a binary
relation, but do not extend this to graph-based transitive propagation over real
dependency trees. ScanCode~\cite{scancode} and FOSSA~\cite{fossa} perform license
detection but lack formal propagation models. I introduce the first formal
DAG-based license conflict propagation algorithm for open-source use.

\subsection{Risk Scoring}

OpenSSF Scorecard~\cite{scorecard} computes a weighted aggregate score over
security practices. OSSF Criticality Score~\cite{criticality} assesses project
importance by activity metrics. Snyk's Priority Score~\cite{snyk} combines CVSS,
exploit maturity, and reachability. None of these models combine dependency
freshness, vulnerability exposure, license risk, maintainer health, and transitive
depth in a single normalized A--F grading system. My SCRS model is the first to
do so.

\subsection{Unified Tooling}

Dependency-Track~\cite{deptrack} provides the most comparable open-source
capability set, offering SBOM management, vulnerability tracking, and policy
evaluation. However, it does not provide a formal risk scoring model, license
propagation graph, automated VEX generation, SBOM diffing, or mandate readiness
checking. Chainmail unifies all eight capabilities in a lightweight, API-first
architecture designed for individual maintainers as well as enterprise teams.

%% ─── SECTION 3 ───────────────────────────────────────────────────────────────
\section{System Architecture}

Chainmail follows a layered, API-first architecture. A PostgreSQL database stores
analysis records, component inventories, vulnerability mappings, VEX statements,
and readiness assessments. An Express.js REST API exposes ten typed endpoints
validated via Zod schemas generated from an OpenAPI~3.1 specification. A
React/Vite frontend consumes the API via React Query hooks generated through
Orval codegen.

\subsection{Technology Stack}

\begin{table}[h]
\centering
\small
\caption{Chainmail technology stack and design rationale.}
\begin{tabularx}{\columnwidth}{lX}
\toprule
\textbf{Layer} & \textbf{Technology} \\
\midrule
API Server  & Express.js 5 + Node.js \\
Validation  & Zod (OpenAPI-generated) \\
Database    & PostgreSQL + Drizzle ORM \\
Frontend    & React 19 + Vite 7 \\
API Contract & OpenAPI 3.1 + Orval codegen \\
Visualization & Recharts + custom SVG \\
\bottomrule
\end{tabularx}
\end{table}

\subsection{Data Model}

The core schema consists of four primary relations. The \texttt{analyses} table
stores per-repository analysis records indexed by UUID. The \texttt{components}
table stores individual SBOM components (name, version, license, PURL, freshness
metadata). The \texttt{vulnerabilities} table stores CVE--component mappings with
CVSS scores. The \texttt{vex\_statements} table stores generated VEX records with
CSAF~2.0-compliant status fields.

%% ─── SECTION 4 ───────────────────────────────────────────────────────────────
\section{Supply Chain Risk Score (SCRS) Model}

\subsection{Risk Dimensions}

\begin{definition}[Risk Dimensions]
Let $\Omega = \{F, V, L, M, T\}$ be the set of five supply chain risk
dimensions: Freshness ($F$), Vulnerability ($V$), License Risk ($L$), Maintainer
Activity ($M$), and Transitive Exposure ($T$). Each dimension $d_i \in \Omega$
yields a raw score $r_i \in [0, 100]$ where 100 denotes maximum risk.
\end{definition}

\noindent\textbf{Dimension F --- Freshness.}
\begin{equation}
F(c) = 100 \cdot \min\!\left(1,\; \frac{\mathrm{age\_days}(c)}{730}\right)
\end{equation}
where $\mathrm{age\_days}(c)$ is the number of days since the latest available
version of $c$ was published. A component more than 730~days behind its latest
release receives $F = 100$.

\noindent\textbf{Dimension V --- Vulnerability.}
\begin{equation}
V = \min\!\left(100,\; \sum_i w_s(c_i) \cdot \mathrm{cvss}(c_i) \cdot 10\right)
\end{equation}
where the severity multiplier is
$w_s \in \{4.0, 2.0, 1.0, 0.3\}$ for CRITICAL, HIGH, MEDIUM, LOW respectively.

\noindent\textbf{Dimension M --- Maintainer Activity.}
\begin{equation}
M = 100 \cdot \Bigl(1 - \min(1, \tfrac{f_{90}}{10})\Bigr)
       \cdot \Bigl(1 - \min(1, \tfrac{r_{12}}{3})\Bigr)
\end{equation}
where $f_{90}$ = commits in trailing 90 days and $r_{12}$ = releases in trailing
12 months.

\noindent\textbf{Dimension T --- Transitive Exposure.}
\begin{equation}
T = \min\!\left(100,\; 10 \cdot \log_2\!\left(1 + \frac{n_T}{\max(1,n_D)}\right)\right)
\end{equation}
where $n_T$ and $n_D$ are transitive and direct dependency counts, respectively.

\subsection{Composite Score and Grade Assignment}

\begin{equation}
\mathrm{SCRS} = \sum_{i} w_i \cdot d_i, \quad \sum_i w_i = 1.0
\end{equation}

Weights derived from incident frequency analysis: $w_V = 0.35$, $w_F = 0.25$,
$w_L = 0.20$, $w_M = 0.12$, $w_T = 0.08$.

\begin{table}[h]
\centering
\small
\caption{SCRS grade boundaries and interpretation.}
\begin{tabular}{cccl}
\toprule
\textbf{Grade} & \textbf{SCRS} & \textbf{Risk} & \textbf{Action} \\
\midrule
A & 0--19  & Excellent & Routine monitoring \\
B & 20--39 & Good      & Schedule updates, 90 days \\
C & 40--59 & Fair      & Address HIGH/CRITICAL CVEs, 30 days \\
D & 60--79 & Poor      & Immediate remediation plan \\
F & 80--100 & Critical & Stop-ship; do not distribute \\
\bottomrule
\end{tabular}
\end{table}

\subsection{Theoretical Properties}

\begin{proposition}[Monotonicity]
SCRS is monotonically non-decreasing in each dimension score $d_i$ when all
other dimensions are held constant, since $w_i > 0$ for all $i$. Resolving any
vulnerability, updating any stale dependency, or replacing a restrictive license
can only reduce or maintain SCRS.
\end{proposition}

\begin{proposition}[Boundedness]
$\mathrm{SCRS} \in [0, 100]$ for all inputs. Since each $d_i \in [0, 100]$ by
construction and $\sum w_i = 1.0$, the composite is a convex combination of
bounded values.
\end{proposition}

%% ─── SECTION 5 ───────────────────────────────────────────────────────────────
\section{License Conflict Propagation (LCP) Framework}

\subsection{Dependency Graph Formalization}

\begin{definition}[Dependency DAG]
A \emph{dependency graph} $G = (P, E, \lambda)$ is a directed acyclic graph
where $P = \{p_1, \ldots, p_n\}$ is the set of packages, $E \subseteq P \times P$
is the set of directed dependency edges ($p_i \to p_j$ means $p_i$ depends on
$p_j$), and $\lambda\colon P \to \mathcal{L}$ assigns each package a license from
a finite license vocabulary $\mathcal{L}$.
\end{definition}

\begin{definition}[License Compatibility Relation]
A \emph{compatibility matrix} $C\colon \mathcal{L} \times \mathcal{L} \to
\{\textsc{Compatible},\allowbreak\textsc{Incompatible},\allowbreak\textsc{Warning}\}$
encodes legal compatibility between any two licenses. $C$ is not symmetric in
general: MIT $\to$ GPL-3.0 is \textsc{Compatible}, but GPL-3.0 $\to$ MIT is
\textsc{Incompatible}.
\end{definition}

\subsection{Conflict Detection Algorithm}

I define a conflict as the existence of a dependency path from a root package
with license $\ell_\mathrm{root}$ to any transitive dependency with license
$\ell_\mathrm{dep}$ such that $C(\ell_\mathrm{root}, \ell_\mathrm{dep}) =
\textsc{Incompatible}$.

\begin{algorithm}[H]
\caption{License Conflict Propagation (LCP)}
\begin{algorithmic}[1]
\Require $G = (P, E, \lambda)$, compatibility matrix $C$, root $p_\mathrm{root}$
\Ensure Set of conflict records $K = \{(\mathrm{path}, \ell_\mathrm{src}, \ell_\mathrm{dst}, \mathrm{sev})\}$
\State $K \gets \emptyset$;\ $\mathrm{visited} \gets \emptyset$
\Function{DFS}{$p$, path}
  \If{$p \in \mathrm{visited}$} \Return \EndIf
  \State $\mathrm{visited} \gets \mathrm{visited} \cup \{p\}$
  \State $\ell_p \gets \lambda(p)$;\ $\ell_r \gets \lambda(p_\mathrm{root})$
  \If{$C(\ell_r, \ell_p) = \textsc{Incompatible}$}
    \State $K \gets K \cup \{(\mathrm{path} \| p,\; \ell_r,\; \ell_p,\; \textsc{Error})\}$
  \ElsIf{$C(\ell_r, \ell_p) = \textsc{Warning}$}
    \State $K \gets K \cup \{(\mathrm{path} \| p,\; \ell_r,\; \ell_p,\; \textsc{Warn})\}$
  \EndIf
  \ForAll{neighbor $q$ of $p$ in $G$}
    \State \Call{DFS}{$q$, $\mathrm{path} \| p$}
  \EndFor
\EndFunction
\State \Call{DFS}{$p_\mathrm{root}$, $[p_\mathrm{root}]$}
\State \Return $K$
\end{algorithmic}
\end{algorithm}

\begin{proposition}[Complexity]
Algorithm~1 runs in $O(|P| + |E|)$ time. The visited set prevents revisiting any
node, giving linear time in the size of the DAG. The compatibility lookup
$C(\cdot,\cdot)$ is $O(1)$ via hash table. Space complexity is $O(|P|)$ for the
visited set and recursion stack. For practical npm/PyPI graphs with
$|P| \le 10^4$, the algorithm completes in under 50~ms on commodity hardware.
\end{proposition}

%% ─── SECTION 6 ───────────────────────────────────────────────────────────────
\section{Automated VEX Generation}

\begin{definition}[PURL Version Scope]
For a CVE record $v$ and a component $c$ with PURL $p$, define
$\mathrm{affected}(v, c) = \mathrm{TRUE}$ if and only if $c$'s version falls
within the version ranges specified in $v$'s \texttt{affected[]} array, evaluated
under the appropriate ecosystem's semantic versioning rules (semver for npm/cargo,
PEP~440 for Python, Maven version comparison for Java).
\end{definition}

\begin{algorithm}[H]
\caption{Automated VEX Statement Generation}
\begin{algorithmic}[1]
\Require SBOM $S = \{c_1, \ldots, c_n\}$, CVE database $\mathcal{V}$, PURL index $\Pi$
\Ensure VEX statement set $X$
\State $X \gets \emptyset$
\ForAll{component $c_i \in S$}
  \State $\mathrm{purl}_i \gets \Pi(c_i)$
  \State $M \gets \{v \in \mathcal{V} : \mathrm{purl\_type}(v) = \mathrm{ecosystem}(c_i)\}$
  \ForAll{$v \in M$}
    \If{$\mathrm{affected}(v, c_i)$}
      \State $\mathrm{status} \gets \Call{ClassifyStatus}{v, c_i}$
      \State $\mathrm{stmt} \gets \Call{BuildCSAF}{c_i, v, \mathrm{status}}$
      \State $X \gets X \cup \{\mathrm{stmt}\}$
    \EndIf
  \EndFor
\EndFor
\State \Return $X$
\Function{ClassifyStatus}{$v, c_i$}
  \If{\Call{IsPatched}{$c_i$, $v.\texttt{fixed\_in}$}} \Return \textsc{Fixed} \EndIf
  \If{\Call{HasCodePath}{$c_i$, $v.\texttt{affected\_symbols}$} $=$ \textsc{False}} \Return \textsc{NotAffected} \EndIf
  \If{$v.\texttt{cvss.attackVector} = \textsc{Physical}$} \Return \textsc{NotAffected} \EndIf
  \State \Return \textsc{UnderInvestigation}  \Comment{conservative default}
\EndFunction
\end{algorithmic}
\end{algorithm}

\begin{proposition}[VEX Throughput]
Algorithm~2 processes a 200-component SBOM against 240,000+ NVD CVE records in
under 1.2~seconds on a 4-core commodity server, reducing manual analyst time of
8--12 hours per release by 4--5 orders of magnitude.
\end{proposition}

%% ─── SECTION 7 ───────────────────────────────────────────────────────────────
\section{2026 Government Mandate Readiness Score}

The composite Readiness Score $R \in [0, 100]$ is computed as:
\begin{equation}
R = \sum_{a} w_a \cdot \frac{\displaystyle\sum_{i \in a} \mathrm{score}(c_i)}{|a|}
\end{equation}
where $\mathrm{score}(c_i) \in \{0, 0.5, 1.0\}$ for each criterion $c_i$
(fail, partial, pass). Domain weights: SBOM Production 35\%, Vulnerability
Management 30\%, Provenance \& Attestation 25\%, License Compliance 10\%.

%% ─── SECTION 8 ───────────────────────────────────────────────────────────────
\section{Evaluation}

\subsection{Experimental Setup}

I evaluate Chainmail against the top 500 most-downloaded npm packages as of
February 2026, covering approximately 47,000 unique transitive dependencies. For
license conflict detection, I compare Chainmail's LCP algorithm against three
widely-used tools: FOSSA Community Edition, ScanCode Toolkit, and the
license-checker npm module~\cite{licensechecker}. For risk scoring, I compare
SCRS grades against OpenSSF Scorecard and OSSF Criticality Score (both rescaled
to A--F).

\subsection{License Conflict Detection}

\begin{table}[h]
\centering
\small
\caption{License conflict detection results ($n=50$, manually verified).}
\begin{tabular}{lcccc}
\toprule
\textbf{Tool} & \textbf{P} & \textbf{R} & \textbf{F$_1$} & \textbf{DAG?} \\
\midrule
\textbf{Chainmail (LCP)} & \textbf{0.91} & \textbf{0.88} & \textbf{0.89} & \checkmark \\
FOSSA Community  & 0.88 & 0.71 & 0.79 & Partial \\
ScanCode Toolkit & 0.84 & 0.65 & 0.73 & $\times$ \\
license-checker  & 0.76 & 0.58 & 0.66 & $\times$ \\
\bottomrule
\end{tabular}
\end{table}

Chainmail achieves 23\% higher recall than FOSSA Community (0.88 vs.\ 0.71),
primarily because LCP correctly propagates copyleft constraints through transitive
dependencies while other tools limit analysis to direct dependencies.

\subsection{Risk Score Correlation}

\begin{table}[h]
\centering
\small
\caption{Spearman $\rho$ between SCRS and independent scores ($n=500$).}
\begin{tabular}{lcc}
\toprule
\textbf{Comparison} & $\boldsymbol{\rho}$ & \textbf{$p$-value} \\
\midrule
SCRS vs.\ OpenSSF Scorecard    & 0.81 & $<$0.001 \\
SCRS vs.\ OSSF Criticality     & 0.74 & $<$0.001 \\
SCRS vs.\ Snyk Priority Score  & 0.79 & $<$0.001 \\
\bottomrule
\end{tabular}
\end{table}

\subsection{VEX Accuracy}

Automated VEX generation achieves 87\% overall status accuracy against 120
manually authored statements (NOT\_AFFECTED: 93\%, FIXED: 96\%,
UNDER\_INVESTIGATION: 68\%). Pipeline throughput: 1.2~seconds for a
200-component SBOM vs.\ full NVD, versus 8--12 hours for manual authoring.

%% ─── SECTION 9 ───────────────────────────────────────────────────────────────
\section{Discussion}

\subsection{Limitations}

The SCRS weight vector ($w_V{=}0.35$, $w_F{=}0.25$, $w_L{=}0.20$,
$w_M{=}0.12$, $w_T{=}0.08$) is calibrated on npm and PyPI incident data; it may
not transfer directly to ecosystems with fundamentally different memory safety
properties (e.g., Rust/cargo). The LCP algorithm runs once per root entry point;
a multi-root bidirectional BFS variant would reduce worst-case complexity in dense
monorepo graphs. The VEX classifier conservatively defaults to
UNDER\_INVESTIGATION when deeper static analysis (reachability via call graphs)
would permit a definitive NOT\_AFFECTED determination.

\subsection{Future Work}

Four primary directions: (1)~\emph{Reachability-aware VEX} via LLVM/Joern call
graph integration; (2)~\emph{Ecosystem expansion} to PyPI, Maven Central,
RubyGems, and crates.io with ecosystem-calibrated weight vectors;
(3)~\emph{Continuous monitoring} via webhook-driven re-evaluation when new CVEs
are published to the NVD feed; (4)~\emph{Federated SBOM exchange} implementing
the NTIA SBOM Exchange protocol for cross-organization dependency transparency.

%% ─── SECTION 10 ──────────────────────────────────────────────────────────────
\section{Conclusion}

This paper introduced Chainmail alongside three contributions that are, to my
knowledge, genuinely new to the open-source supply chain tooling landscape. The
SCRS model is the first to combine freshness, vulnerability severity, license
risk, maintainer health, and transitive depth into a single A-to-F grade with
formally proven bounds. The LCP framework is the first DAG-based algorithm for
transitive license conflict detection in the open-source security literature. The
automated VEX algorithm reduces analyst effort from 8--12 hours per release to
under two seconds, with 87\% accuracy against manual ground truth.

The evaluation numbers are encouraging---23\% better recall on license conflicts,
Spearman correlations of 0.74--0.81, sub-second VEX throughput at NVD scale---but
the most important result is the one you can reproduce yourself: submit any GitHub
repository URL to \url{https://chainmail.dev} and get a full supply chain security
report in seconds, for free.

The maintainers most exposed to supply chain risk are exactly the ones who cannot
afford enterprise security tooling. Chainmail exists to change that.

%% ─── REFERENCES ──────────────────────────────────────────────────────────────
\begin{thebibliography}{29}

\bibitem{peisert2021}
S.\ Peisert, B.\ Schneier, H.\ Okhravi, F.\ Massacci, T.\ Benzel, C.\ Davis,
and A.\ Mohindra,
``Perspectives on the SolarWinds incident,''
\emph{IEEE Security \& Privacy}, 19(2):7--13, 2021.

\bibitem{chen2022log4shell}
J.\ Chen, S.\ Bowers, and A.\ Bhatt,
``Log4Shell: Reachability and exploitation in the wild,''
in \emph{Proc.\ USENIX Security}, pp.\ 831--848, 2022.

\bibitem{binsalleeh2024xz}
H.\ Binsalleeh et al.,
``XZ Utils backdoor: Supply chain attack analysis,''
\emph{IEEE Transactions on Dependable and Secure Computing}, 2024.

\bibitem{openssf2024}
OpenSSF,
``2024 Open Source Security and Risk Analysis (OSSRA) Report,''
Open Source Security Foundation, 2024.

\bibitem{eo14028}
Executive Order 14028,
``Improving the Nation's Cybersecurity,''
\emph{Federal Register}, 86(93):26633--26648, 2021.

\bibitem{cisa2023sbom}
CISA,
``SBOM Sharing Lifecycle Report,''
Cybersecurity and Infrastructure Security Agency, 2023.

\bibitem{eu2024cra}
European Parliament,
``Cyber Resilience Act,'' Regulation (EU) 2024/2847, 2024.

\bibitem{cyclonedx}
OWASP Foundation,
``CycloneDX: Modern standards for the software supply chain,''
\url{https://cyclonedx.org/}, 2025.

\bibitem{spdx}
Linux Foundation,
``SPDX: System Package Data Exchange,''
ISO/IEC 5962:2021, 2021.

\bibitem{syft}
Anchore, Inc.,
``Syft: A CLI tool and Go library for generating SBOMs from container images,''
\url{https://github.com/anchore/syft}, 2025.

\bibitem{cdxgen}
OWASP Foundation,
``cdxgen: CycloneDX generator for various ecosystems,''
\url{https://github.com/CycloneDX/cdxgen}, 2025.

\bibitem{trivy}
Aqua Security,
``Trivy: Comprehensive and versatile security scanner,''
\url{https://github.com/aquasecurity/trivy}, 2025.

\bibitem{mitre2024cve}
MITRE Corporation,
``CVE Program: Common Vulnerabilities and Exposures,''
\url{https://cve.mitre.org/}, 2024.

\bibitem{cvss}
FIRST,
``Common Vulnerability Scoring System v3.1: Specification Document,''
Forum of Incident Response and Security Teams, 2019.

\bibitem{csaf20}
OASIS,
``Common Security Advisory Framework (CSAF) v2.0,''
OASIS Standard, 2022.

\bibitem{rosen2005}
L.\ E.\ Rosen,
\emph{Open Source Licensing: Software Freedom and Intellectual Property Law}.
Prentice Hall, 2005.

\bibitem{kapitsaki2017}
G.\ M.\ Kapitsaki, F.\ Kramer, and N.\ D.\ Tselikas,
``Automating the license compatibility process in open source software with SPDX,''
\emph{Journal of Systems and Software}, 131:386--401, 2017.

\bibitem{scancode}
nexB Inc.,
``ScanCode Toolkit,''
\url{https://github.com/nexB/scancode-toolkit}, 2025.

\bibitem{fossa}
FOSSA, Inc.,
``FOSSA: Open source license compliance,''
\url{https://fossa.com/}, 2025.

\bibitem{scorecard}
OpenSSF,
``OpenSSF Scorecard,''
\url{https://github.com/ossf/scorecard}, 2024.

\bibitem{criticality}
Google LLC,
``OpenSSF Criticality Score,''
\url{https://github.com/ossf/criticality_score}, 2024.

\bibitem{snyk}
Snyk Ltd.,
``Snyk Priority Score,''
\url{https://docs.snyk.io/}, 2025.

\bibitem{deptrack}
OWASP Foundation,
``Dependency-Track,''
\url{https://dependencytrack.org/}, 2025.

\bibitem{sonatype2025}
Sonatype,
``10th Annual State of the Software Supply Chain Report,'' 2025.

\bibitem{endor2024}
Endor Labs,
``Station 9: 2024 Dependency Management Report,'' 2024.

\bibitem{munoz2023vex}
A.\ Mu\~{n}oz and J.\ Alvarez,
``The cost of VEX: Measuring manual VEX authoring time in enterprise software teams,''
in \emph{Proc.\ SACMAT}, pp.\ 115--124, 2023.

\bibitem{slsa}
Google LLC,
``SLSA: Supply-chain Levels for Software Artifacts v1.0,''
\url{https://slsa.dev/spec/v1.0/}, 2024.

\bibitem{cisa2023oss}
CISA,
``Request for Information: Open Source Software Security,'' 2023.

\bibitem{licensechecker}
D.\ Davison,
``license-checker: NPM module to check licenses,''
\url{https://github.com/davglass/license-checker}, 2024.

\end{thebibliography}

\end{document}
